\section{A comparison scale carrying diffuse Weyl curvature}
\label{sec:spectral-windows}

The median construction in Theorem~\ref{thm:band-charge} controls
the entire modified volume when its spectral margin is positive.
For descent it suffices to find a single region carrying a fixed
positive chiral mass. We now obtain one throughout the strict band.

\begin{theorem}\label{thm:spectral-windows}
Let \(g_i\) be unit-volume positive Yamabe metrics on a fixed
nonstandard smooth oriented homotopy four-sphere with
\(e_i\to I=I^+(M)<2T\). After passing to a subsequence let
\(Y_i\to y>0\). There are constant rescalings \(k_i\) of smooth
positive-Ricci metrics in these classes and basepoints for which
the pointed limit \(X\) is noncollapsed. Write \(m=\mathcal H^4\)
in Euclidean volume normalization. In this comparison let
\(\sigma,\lambda^+,\lambda^-\) be the local weak limits of
\(L^2dV,|W^+|^2_{\rm HS}dV,|W^-|^2_{\rm HS}dV\) on the sources.
The absolutely continuous parts relative to \(m\) satisfy
\[
 \sigma_{\rm ac}(X)\ge y^2/36,\qquad
 (\lambda^+_{\rm ac}+\lambda^-_{\rm ac})(X)\ge y^2/24 .
\]
For every \(\epsilon,\zeta,r_*>0\) a fixed finite number of
smoothly parametrized disjoint standard disks on each late source,
with union \(U_i\), can be chosen in comparison balls of fixed radii
\(r_j<r_*\), each having original full curvature integral below
\(\zeta\). One sign \(s\), fixed along the subsequence, satisfies
\[
 \lambda_i^s(U_i)\ge Y_i^2/48-\epsilon
                  \ge 2T-e_i-\epsilon.
\]
On the same disks both signs satisfy
\(\lambda_i^\pm(M\setminus U_i)\le 2e_i-Y_i^2/24+\epsilon\).
The second bound is informative for each sign when \(Y_i^2/24>e_i\).
\end{theorem}

\begin{proof}
Use the volume-one modified representatives \(q_i\), with
\(\mu_i=dV_{q_i}\), \(R_{q_i}=6L_{q_i}+c_i\), \(c_i\le0\),
from Section~\ref{sec:endpoint-charge}. The positive-Ricci comparisons
exist by Section~\ref{sec:physical-volume}. Put
\(\sigma_i=L_{q_i}^2\mu_i\), a conformal measure of bounded total
mass. Discard finitely many indices and fix
\(0<6\sqrt\beta<Y_0\le Y_i\) using the strict subcritical budget.
Start with the function \(h_i=1\). At each step it is Lipschitz,
\(0\le h_i\le1\), \(\int h_i^4d\mu_i\to1\), and
\(\int|dh_i|^4dV=o(1)\), the latter integral being conformally
invariant. The restricted mass \(\sigma_i(\{h_i>0\})\) must
exceed \(\beta\) eventually: otherwise the ordinary Yamabe
inequality tested on \(h_i\) gives
\[
 (Y_0-6\sqrt\beta)\|h_i\|_{4,q_i}^2
   \le 6\|dh_i\|_{2,q_i}^2=o(1)
\]
on a subsequence. Here and below we use \(\mu_i(M)=1\).

Scale the positive-Ricci comparison so that the least radius
enclosing \(\beta\) of this restricted measure is one, witnessed
at \(p_i\). The restricted measure is absolutely continuous on
each smooth source, so continuity of ball mass and compactness of
the center set give the normalization. If some bounded centered
window carries a positive lower bound of \(h_i^4\mu_i\) on a
subsequence, keep this scale. Otherwise every fixed window has
weighted volume tending to zero. Choose \(R_i\uparrow\infty\)
sufficiently slowly that \(\int_{B_{R_i^2}(p_i)}h_i^4d\mu_i=o(1)\).
Multiply \(h_i\) by a distance cutoff zero through \(R_i\), rising
linearly in log radius to one at \(R_i^2\). Its fourth-power
gradient integral is at most \(2\pi^2/(\log R_i)^3\) by the
Bishop area bound. The new functions preserve both limits and
their positive sets lose at least \(\beta\) of spectral mass.
Uniform boundedness of \(\sigma_i(M)\) forces termination after
finitely many such deletions.

On the retained scale put \(\theta_i=h_i^4\mu_i\). For a Lipschitz
cutoff \(0\le\chi\le1\) in a unit ball the Yamabe test gives
\[
 (Y_0-6\sqrt\beta)\Big(\int\chi^4d\theta_i\Big)^{1/2}
       \le6\|d(h_i\chi)\|_{2,q_i}^2 .
\]
The norm on the right before squaring is at most
\(o(1)+\|d\chi\|_{4,k_i}\). Inward logarithmic cutoffs give
a uniform upper bound tending to zero for limiting mass on balls
of radius tending to zero. On the other hand a fixed bounded
window carries positive weighted mass. The relative Bishop packing
bound covers it by boundedly many quarter-balls. A linear cutoff on
the half-ball about a quarter-ball of definite mass now gives a
positive lower bound on the comparison volume of that half-ball.
Relative volume comparison gives noncollapse on all bounded
windows. Pass to a pointed limit and common local weak limits of
the measures, including \(\theta_i\), the unrestricted \(\mu_i\)
and the original full curvature measures \(\nu_i\). The preceding
estimates give \(v:=\theta(X)>0\) and \(\theta(\{x\})=0\).

Off the countable atoms of the unrestricted \(\sigma\), choose a
small continuity ball at each point. For \(q_i=w_i^2 k_i\),
\(-\Delta_{k_i}w_i\le L_{q_i}w_i^3\), and the squared local
\(L^2(k_i)\) norm of the potential \(L_{q_i}w_i^2\) is the
original spectral mass on the ball. The \(p=4\) estimate
\eqref{eq:physical-moser} therefore absorbs this term on a sufficiently
small fixed ball and bounds \(w_i^4\) in \(L^2(k_i)\) on its half.
Noncollapsed volume convergence and Cauchy--Schwarz show that the
limit of \(\mu_i\) is absolutely continuous there. Since
\(\theta_i\le\mu_i\), countable covering and the absence of
\(\theta\)-atoms give \(\theta\ll m\) on all of \(X\).
This uses only local convergence.

Set \(\alpha_i=L_{q_i} h_i^2\mu_i\) and take its local weak
limit \(\alpha\). The matrix of measures with diagonal
\((\theta_i,\sigma_i)\) and off-diagonal \(\alpha_i\) is
positive semidefinite: every constant quadratic test has a square
as its density. The property passes to local weak limits. Taking
densities relative to \(\theta+\sigma\) (Cauchy--Schwarz also
gives \(\alpha\ll\theta+\sigma\)) and testing rational pairs yields
\(\alpha=f\theta\) and \(\sigma\ge f^2\theta\).
For \(R>1\), use the centered cutoff \(\chi_R\) equal to one
through \(R\), falling linearly in log radius to zero at \(R^2\).
Testing the Yamabe inequality on \(h_i\chi_R\) gives
\[
 Y_i\Big(\int\chi_R^4d\theta_i\Big)^{1/2}
 \le 6\|d(h_i\chi_R)\|_{2,q_i}^2 +6\int\chi_R^2d\alpha_i .
\]
At fixed \(R\) the radial tests converge on the common pointed
realization. The gradient norm is at most
\(o_i(1)+(2\pi^2)^{1/4}(\log R)^{-3/4}\).
Taking the source limit and then \(R\to\infty\) gives
\(y\sqrt v\le6\alpha(X)\). Cauchy--Schwarz and \(\theta\ll m\)
now give \(\sigma_{\rm ac}(X)\ge\int f^2d\theta
\ge\alpha(X)^2/v\ge y^2/36\). The pointwise algebraic inequality
\(\lambda_i^++\lambda_i^-\ge(3/2)\sigma_i\) passes to measures
and to their absolutely continuous parts.

Take \(s\) maximizing \(\lambda^s_{\rm ac}(X)\).
Both absolutely continuous parts give full measure to regular
points outside atoms of \(\nu\). Apply the finite ball selection
in Section~\ref{sec:physical-volume} to their sum together with
\(\sigma_{\rm ac}\) on a bounded
continuity window containing all but a prescribed tolerance of its
mass. Use differentiation of this density for the loss on shrinking
the balls. The same loss estimate bounds the omitted part of
either absolutely continuous measure. On the sources use the
smooth embedded disks containing the fixed inner comparison balls
and contained in the fixed outer ones. Lower bounds transfer by
weak convergence of the full chiral measures at common continuity
radii; the curvature ceiling uses \(\nu_i\) of \(g_i\).
Subtracting the opposite total mass \(e_i\) from the combined
lower bound proves the two individual residual estimates. Adjust
the tolerance using \(Y_i\to y\) and the budget
\(Y_i^2\ge48(2T-e_i)\).
\end{proof}

There is a second estimate for the full residual. Take the same fixed
disks using the absolutely continuous spectral measure as above.
With \(\Delta_i=36\ell_i^2-Y_i^2\),
\[
   \sigma_i(M\setminus U_i)\le \Delta_i/36+\epsilon,\qquad
   \lambda_i^\pm(M\setminus U_i)\le \Delta_i/6+6\epsilon .
\]
Indeed the transfer captures at least \(Y_i^2/36-\epsilon\) of
\(\sigma_i\), and each Weyl block satisfies \(|W^\pm|^2\le6L^2\).
In particular if \(\Delta_i\to0\), both chiral measures can be
captured to arbitrary tolerance in this one frame. This controls
escape, singular mass and the excluded atoms together. The selected
finite disks and their radii are fixed before the source limit.

For the selected sign, escape together with the entire singular part
in this frame costs at most \(I-y^2/48\). Its raw complement on the
finite family is at most \(2(e_i-T)+\epsilon\), hence is strictly
below \(T\) when \(I<3T/2\) and the tolerance is small. The chiral
density need not be positive in both signs at the same regular point.
Taking a curvature ceiling \(\zeta<2I\) makes every individual disk
carry less than \(I/2\) in either sign. At fixed tolerance the number
of disks is fixed, so a further subsequence has one label carrying a
fixed positive mass of sign \(s\), with a fixed positive mass of the
same sign on its other side. This marked sphere is a candidate for
one repair on the original type.

\begin{corollary}[A marked bicollar with small incoming curvature]
\label{cor:thin-raw-collar}
For every \(\eta,\zeta,r_*>0\) the selected sign admits a fixed
comparison radius \(r<r_*\) and a number \(b>0\) such that on all
late sources a parametrized disk \(D_i\) and a two-sided smooth
bicollar \(A_i\) of its boundary lie in a single comparison ball
\(B_i(r)\). Both complete sides carry at least \(b\) of sign \(s\),
and, writing \(\nu_i=|\Rm_{g_i}|_{\rm tens}^2dV_{g_i}\),
\[
 \nu_i(B_i(r))<\zeta,\qquad
 (\nu_i+\lambda_i^s)(A_i)<\eta\,\lambda_i^s(D_i).
\]
\end{corollary}
\begin{proof}
Choose a regular point of differentiation of \(m,\lambda^s,\nu\)
where \(d\lambda^s_{\rm ac}/dm\) is finite and positive and
\(d\nu_{\rm ac}/dm\) is finite. There the singular parts have zero
four-dimensional density. Choose \(0<\tau<1/6\) so the limiting
density ratio of the mass of \(\nu+\lambda^s\) on
\(B_r\setminus B_{(1-3\tau)r}\) to that of \(\lambda^s\) on
\(B_{(1-3\tau)r}\) is below \(\eta/2\). Then choose a sufficiently
small fixed joint continuity radius meeting the comparison volume
threshold at \((1-\tau)r\) and a strict outer curvature ceiling
below \(\min(\zeta,2I)\). The finite-radius transfers and smooth
open model ball of Section~\ref{sec:physical-volume} apply on
every late source. Its interior contains the closed inner ball
compactly. In its parameter domain choose three nested compact
concentric disks, the smallest containing that compact preimage
in its interior. Use the middle disk for \(D_i\) and the annulus
between the other two for its bicollar. The displayed estimate
follows by inclusion. The fixed inner ball supplies a positive
lower mass, while the entire outer ball carries less than \(I/2\)
of sign \(s\). Taking \(b\) smaller than both margins proves the claim.
\end{proof}

The marking is smooth on both faces. Its normalized distortion and
differentiated coordinate residual can depend on its actual width.
Even a collar of arbitrarily small raw curvature can require a large
change in these quantities. If the entire protected core lies in
the image of one fixed compact of a smooth locally conformally flat
comparison satisfying Theorem~\ref{prop:local-min}, strong
local chiral convergence would give core energy tending to zero.
This excludes a round restart covering the charged core in such
a marking. The evacuated graft of Proposition~\ref{prop:graft}
can still occur on another region.
